% ================================================================
% Section: Taxonomy of Small Language Models
% ================================================================

\section{A Faceted Taxonomy of Small Language Models}
\label{sec:taxonomy}

\subsection{The Scheme}
\label{subsec:taxonomy_scheme}

We index the \gls{slm} design space along four \emph{primary} facets, architecture, training strategy, modality, and application domain (Figure~\ref{fig:taxonomy}), and five \emph{secondary} facets that govern adoption: licensing, data transparency, safety alignment, reasoning mode, and deployment target (Table~\ref{tab:taxonomy_facets} in \suppsec~\ref{supp:tables}). The facets are \emph{separately observable}, in that each can be read off a model's primary source, but they are \emph{not orthogonal}: pruning-plus-distillation presupposes a dense teacher, and a vision encoder ties modality to architecture. The scheme is also \emph{multi-label}: Gemma~3 is distilled \emph{and} trained on a curated corpus, Granite~4.0 is a hybrid state-space model \emph{and} a mixture of experts. It is an indexing scheme illustrated with exemplars, not a partition of the inventory. The released inventory records per variant the fields its sources state directly (modality, access, license, dense or mixture-of-experts, post-training status, context length), and the figures are drawn from those fields; architecture and training facets are assigned in this paper to representative models only. The application-domain facet is the weakest of the four, since it usually describes use rather than the artifact, so we assign it only when the developer trained the released artifact for a domain and says so.

\begin{figure}[t]
\centering
\footnotesize
% Palette shared with the generated figures (source/scripts/generate_figures_v2.py).
\definecolor{txInk}{HTML}{34404A}
\definecolor{txSlate}{HTML}{A2AAB1}
\definecolor{txIndigo}{HTML}{7397B8}
\definecolor{txOcean}{HTML}{D8A879}
\definecolor{txTeal}{HTML}{79A88E}
\definecolor{txCoral}{HTML}{9B8FB5}
\newcommand{\txax}[2]{{\normalsize\bfseries\color{#1!88!black}#2}}
\newcommand{\txcat}[1]{{\scriptsize\bfseries\color{txInk}#1}}
\newcommand{\txsep}{{\color{txSlate}\,\textendash\,}}
\newcommand{\txex}[1]{{\scriptsize\itshape\color{txSlate}#1}}
\begin{tikzpicture}[
  panel/.style={draw=#1!55, line width=0.7pt, rounded corners=3pt, inner sep=5pt,
                text width=0.455\linewidth, align=left, fill=#1!7},
  node distance=3mm and 3mm,
]
\node[panel=txIndigo] (arch) {%
  \txax{txIndigo}{1.\ Architecture}\\[2pt]
  \txcat{Dense transformer}\txsep GQA, RoPE, gated FFN\\ \txex{Phi-4-mini, Qwen3, SmolLM3, Gemma 3}\\[1pt]
  \txcat{Mixture-of-experts}\txsep sparse activation\\ \txex{Granite 4.0 Tiny; small-active MoE: Mixtral, Qwen3-30B-A3B}\\[1pt]
  \txcat{Hybrid SSM/attention}\txsep linear-time\\ \txex{Falcon-H1, Granite 4.0, Nemotron Nano 2}%
};
\node[panel=txOcean, right=of arch] (train) {%
  \txax{txOcean}{2.\ Training strategy}\\[2pt]
  \txcat{From scratch}\txsep curated / synthetic data\\ \txex{Phi, Qwen3, SmolLM3}\\[1pt]
  \txcat{Distillation}\txsep transfer from a larger teacher\\ \txex{Gemma 3/4, DeepSeek-R1-Distill, Phi-4-mini-reasoning}\\[1pt]
  \txcat{Pruning + distillation}\txsep shrink a larger model\\ \txex{Minitron, Sheared-LLaMA, Llama 3.2 1B/3B}%
};
\node[panel=txTeal, below=of arch] (mod) {%
  \txax{txTeal}{3.\ Modality}\\[2pt]
  \txcat{Text-only}\txsep majority of SLMs\\ \txex{Mistral 7B, SmolLM3, Qwen3}\\[1pt]
  \txcat{Vision--language}\txsep image encoder added\\ \txex{Phi-3-vision, Gemma 3 (4B+), PaliGemma}\\[1pt]
  \txcat{Text + vision + audio}\txsep one sub-10B model\\ \txex{Gemma 3n, Gemma 4, Voxtral Mini (audio)}%
};
\node[panel=txCoral, below=of train] (dom) {%
  \txax{txCoral}{4.\ Application domain}\\[2pt]
  \txcat{General-purpose}\txsep broad benchmarks\\ \txex{Phi, Gemma, Llama, Qwen, Mistral}\\[1pt]
  \txcat{Healthcare}\txsep clinical text\\ \txex{Meerkat, OptimCLM (encoders)}\\[1pt]
  \txcat{Code / multilingual}\txsep domain corpora\\ \txex{StarCoder2, DeepSeek-Coder, Aya 23}%
};
\end{tikzpicture}
\caption{The four primary facets of the taxonomy. Facets are separately observable but not orthogonal, and assignment is multi-label. Representative models are italicised; \suppsec~\ref{supp:taxonomy} develops each facet and Table~\ref{tab:taxonomy_summary} gives the tabular form.}
\Description{A two-by-two arrangement of four labelled panels, one per primary taxonomy facet: architecture, training strategy, modality, and application domain. Each panel lists three categories with a short defining characteristic and representative models.}
\label{fig:taxonomy}
\end{figure}

% ----------------------------------------------------------------
\subsection{Architecture}
\label{subsec:taxonomy_architecture}

\emph{Dense transformers} remain the majority, from SmolLM2-135M \citep{allal2025smollm2} to 7--14B models. The recurring refinements are grouped-query attention \citep{ainslie2023gqa}, which shrinks the \gls{kv} cache by the ratio of query to \gls{kv} heads; rotary position embeddings with an enlarged base or rescaling for long context; and gated feed-forward layers (SwiGLU in most families, GeGLU in Gemma). Among \glspl{slm} whose records state a context length, the median grew from 2K tokens in 2023 ($n=5$) to 16K in 2024 ($n=39$) and 128K in 2025 ($n=46$) and 2026 ($n=20$); coverage is thin for 2023, so the starting point is uncertain (\suppsec~\ref{supp:figures}).

\emph{Mixtures of experts} \citep{shazeer2017moe,fedus2022switch} decouple total from active parameters, and they sit on the fault line that C2 marks. Qwen3-30B-A3B \citep{qwen3_2025} costs about what a 3B dense model costs per token but needs 17--22~GB of 4-bit weight storage, so it is a small-active-parameter \gls{moe} in this paper, not an \gls{slm}; Granite~4.0 Tiny (about 7B total, 1B active) \citep{granite4_2025} satisfies both conditions. The 2026 releases concentrated in this class, with Nemotron~3.5 Lightning (30B/3B), Laguna XS~2.1 (33B/3B), Ling~3.0 Flash (124B/5.1B), and DeepSeek V4 Flash (284B/13B) appearing between June and August 2026.

\emph{Hybrid state-space models} combine attention with linear-time state-space layers \citep{gu2023mamba}. Jamba \citep{lieber2024jamba} interleaves Mamba and attention blocks and reports a \gls{kv} cache eight times smaller than Mixtral's at 256K context; Falcon-H1 \citep{falconh1_2025}, Granite~4.0 \citep{granite4_2025}, Nemotron Nano~2 \citep{nvidia2025nemotron_nano2}, and the Qwen3.5 small series \citep{qwen3_5_2026} brought the design to production at small scale in 2025--2026. \suppsec~\ref{supp:taxonomy} gives the architectural detail for each class.

\subsection{Training Strategy, Modality, and Domain}
\label{subsec:taxonomy_training}
\label{subsec:taxonomy_modality}
\label{subsec:taxonomy_domain}

Three training routes recur and are often combined: training from scratch on curated or synthetic data (Phi, SmolLM, Qwen), distillation from a larger teacher (Gemma, the DeepSeek-R1 students, Phi-4-mini-reasoning), and pruning a larger model followed by distillation (Minitron, Sheared-LLaMA, Llama~3.2 1B and 3B). Text-only models still dominate the inventory, vision--language models are common at 2--12B, and audio and natively multimodal mobile models (Qwen2.5-Omni, Gemma~3n and~4, Voxtral Mini) grew in 2025--2026 (\suppsec~\ref{supp:figures}). Most \glspl{slm} are released as general-purpose models; developer-trained domain models concentrate in code, healthcare, and multilingual use. \suppsec~\ref{supp:taxonomy} lists representative models for each category.
